\section{家户怎样穿过帝国：田、役、迁徙与信仰}
\label{sec:han-households-mobility}

汉帝国的年代并不只写在皇帝改元和将军出塞的地方。它也写在一户人家能否保住
耕作的人手，能否把粮食送到市上，能否从战区离开，又能否在新的地方被承认为
可居住、可纳役的人。国家通过户口、田租、算赋和徭役来认识这种生活，地方社会
则通过亲属、邻里、债务、雇佣、佃作和保护关系来应付它。二者并非各自独立：
一张簿籍要靠乡里和吏员填写，一次征发要靠道路和家户支撑，一项赈贷或赦令也要
先被人相信能够到达。

\subsection{簿籍看见的是谁}

前二年元始户口数字和建武十五年的度田，都给人一种国家已把人口与土地尽收在
眼前的印象。实际上，两者首先显示的是朝廷以郡国簿籍申报自己所能看见的人户；
它们并不等于真实人口、财富或家庭负担的透明照片。尹湾文书、居延汉简及其他
局部遗存让人看到，分类、出入、转递和更正本身就是耗费劳力的行政工作。某一
县写得清楚，不代表所有县都以同一口径、同一时点完成登记；国家对人户的可见性
恰恰由这些不同的工作速度构成。

这也解释了为何度田会引起争执。清查垦田不是对一片中性土地丈量面积而已，它
会触及谁能以何种身份占有土地，谁被算作家内劳力，谁在战争和迁徙后仍能留在
旧籍，谁必须接受新的地方保护。东汉限制战乱中掠卖者继续为奴婢的诏令同样如此：
它留下朝廷试图承认某类申诉的文字，不能量出获释人数或保证每一处县廷都愿意
推翻已经形成的占有关系。土地、人身与文书在这里交错，不能把“复兴”只读成
一项军事胜利。

\subsection{钱、盐和一条未被看见的运输线}

五铢钱、盐铁、均输和平准经常在制度史中被并列为武帝的财政措施。从读者的视角
看，它们首先是一组移动物品的安排：铜料与钱范要抵达铸处，盐铁要经过产地与
关卡，粮食和其他货物要有人估价、转运和入仓。国家因此获得了较强的调度能力，
也让关吏、商人、手工业者、车船与消费者处在更密集的检查和价格关系中。《史记
·平准书》《汉书·食货志》与《盐铁论》能保存政策语言和争辩，钱范、窖藏与简牍
能显示局部的物质痕迹；没有一种材料能把全国市场规模或每一家户的得失精确化。

这份限度并不使财政史变得无关紧要。恰好相反，它让读者知道远征的成本并不止于
战场。河西一座塞城要粮、马、薪和文书，关中要为它筹钱，沿路的中介要决定是否
按时传递。轮台诏拒绝继续某些远距屯田和出击，显示朝廷至少公开承认这种链条
有尽头；东汉凉州和并州的反复军费，则显示若没有稳定的共同支付办法，地方能够
筹粮、招人和保护道路的网络会得到更大的议价位置。财政集中并非国家越来越强的
单向故事，而是一种把风险重新分配给不同地方和不同人群的过程。

\subsection{迁徙不是战争的余波}

西汉迁徙关东豪族入关中、河西安置降附部众、岭南与西南置郡后的人员移动，都
说明政治中心会主动改变人们的去处。诏令常保存迁徙的目的和名目，却很少留下
被迁者如何安顿田土、亲属和职业的连续声音。战乱所造成的流动更难如此记录：
楚汉战争、王莽末年的崩解、交趾战事、羌乱和汉末的军阀战争，都使逃避征发、
寻找粮食和投靠保护者成为家户的现实选择。不能从某一次“受降”或“平定”推断
流散者已经回到原籍；也不能从史书缺少姓名便假定他们只是在国家叙事中被动移动。

迁徙改变了地方政治。新来者需要租佃、婚姻、市场和安全，原有的家族和官吏需要
决定谁可以进入编户、谁可以使用水源和道路。豪强的作用正应放在这里理解：他们
可能为官府提供消息、出资或维持秩序，也可能将土地、佃客与武装资源集中到不易
核验的网络中。把他们统称为“反国家”或“国家代理人”都会抹掉这种双重性。到
东汉末年，州郡与地方势力之所以能各自筹兵，并不只是野心变多，而是战争、迁徙
和供给已让保护关系先于统一命令分散。

\subsection{祖先、经书、治病与结社}

汉人的宗教和知识生活也不能从朝廷与社会的二分法理解。祭祀、墓葬、家族记忆、
黄老方术、经学、地方捐献、早期佛教供养和治病仪式，往往经过同一批有土地、
书写能力或地方声望的人。楚王英案中的佛教供养语汇，是官方文字捕捉到的一条
早期线索；它能提示某些物品、称谓和宫廷疑虑，不能证明全国已有一套统一教团。
石经残片同样只能证明刻立和展示，不能告诉我们每个读书人如何解释经典。

太平道在一八四年前后成为政治问题，正因符号、治病、师徒关系和末世语言可以
进入已经承受赋役、灾荒和保护压力的乡里。黄巾起事并不证明所有参与者共享同一
信仰，更不允许把所有宗教结社倒写成叛乱的预备队。它提示的只是：当常规的申诉、
救济和保护关系失去可信度时，人们会在官署之外寻找能组织互助与行动的语言。
汉末的地方军事化要放在这条社会史线上读，才不会只剩下几位将领的移动。

\subsection{妇女、儿童与材料的静默}

家户并不是只由成年男性与国家相接。婚姻、继承、丧服、劳作、买卖和战乱中的
离散，都使妇女、儿童与年长者参与资源的分配；然而帝纪、军功传和行政簿籍极少
连续保存他们自己的叙述。吕后、窦太后和邓太后在宫廷材料中格外可见，是因为
她们处在继统与诏令的中心，不能由此推定普通妇女同样能够进入政治文书。相反，
限制掠卖奴婢的诏令、墓志与局部简牍只能提醒我们，战乱如何改变依附和家庭关系，
却不足以算出受影响者的人数，或替她们安排统一的感受。

承认这种静默不是把社会史撤回抽象的“家庭”。它要求读者在看到一次迁徙、一次
征役或一次赦令时追问：谁必须离开土地，谁留在原处维持生产，谁有资格向县廷
提出申诉，谁的名字没有进入记录。答案常常无法完整取得，但问题能够防止帝国
史只保留那些已经拥有军队、印绶或笔墨的人。两汉反复重建簿籍和命令，正是因为
国家从未能无代价地把这些不均等的生活纳入同一张表。

\begin{causalchain}[一户人家的压力怎样进入帝国史]
田、户和货物被登记，才能成为税役、军粮与赈贷的对象；登记需要乡里、文吏、
道路和地方保护才能成立；战争、灾荒与迁徙又不断改变谁仍在原处、谁愿意合作。
当中枢无法稳定维护这些条件，筹粮、收留、传信和募兵的地方网络便会获得新的
政治分量。这个反馈回路解释了两汉多次“恢复”为什么总要回到人户、道路与中介，
而不能由一个皇帝登位或一次战役终结。
\end{causalchain}

\begin{textwindow}[社会材料留下的不是全景照片]
《食货志》、诏令和户口申报说明朝廷希望怎样分类家户；简牍、钱币、碑刻、墓葬
与宗教文本让某些岗位、地点和物件重新可见。它们的尺度不相同：一片简牍不能
换算帝国人口，一个户口数字也不能替家庭发言。把材料并读的目的，是看见不同
尺度相互限制，而不是把沉默填满。
\end{textwindow}
